% !TeX root = ../geometry_main.tex
%==========================================================
%=========================================================
\chapter{Fundamentals of differential Geometry}
%=========================================================
%=========================================================


In this long chapter, we collect various random topics in differential geometry that I feel like writing down.


%-----------------------------------------------------------------------
%========================================================
\section{Basics} \todotag{To be written}
%=========================================================
%----------------------------------------------------------------------


%========================================================
\section{Smoothness \& Whitney embedding theorem} \todotag{To be written}
%=========================================================



%-----------------------------------------------------------------------
%========================================================
\section{Sub-bundles, flows \& Frobenius theorem} \todotag{To be written}
%=========================================================
%-----------------------------------------------------------------------


\begin{mytheorem}[Commutation of vector fields, flows \& straightening coordinates]
The offset of the flow of two vector fields $X,Y \in \mathfrak{X}(M)$ reads in a chart
\begin{align}
\begin{aligned}
    \phi_{X}^{\epsilon} \circ \phi_{Y}^{\delta} - \phi_{Y}^{\delta} \circ \phi_{X}^{\epsilon}= \epsilon \delta [X,Y] + \mathcal{O}(\epsilon^2, \delta^2).
\end{aligned}
\end{align}
The following conditions are equivalent for the commutation of flows
\begin{itemize}
    \item The flows commute
        \begin{align}
            \phi_{X_1}^{\epsilon_1}\circ\dots\circ \phi_{X_k}^{\epsilon_k} = \phi_{X_{\sigma_1}}^{\epsilon_{\sigma_1}}\circ\dots\circ \phi_{X_{\sigma_k}}^{\epsilon_{\sigma_k}} \quad \forall \epsilon_i \in \mathbb{R},\,\, \forall \sigma \in \mathcal{S}_k.
        \end{align}
    \item The vector fields commute, i.e. all Lie brackets vanish
        \begin{align}
            [X_i, X_j] = 0 \quad \forall i,j=1\dots k.
        \end{align} 
    \item There exist local coordinates $\{x^1, \dots, x^n\}$ such that
        \begin{align}
            X_i = \partial_{x^i} \quad \forall i=1\dots k.
        \end{align}
\end{itemize}
\end{mytheorem}
\begin{proof}
The first equation is follows from Taylor expanding the flows.
Point 3 $\Rightarrow$ 1, 2 trivially follows from the definitions.
Point 1 $\Rightarrow$ 2 follows from the first equation by taking the limit $\epsilon, \delta \to 0$.
For point 2 $\Rightarrow$ 1 take a chart where $X_1 = \partial_{x^1}$, then $\phi_{X_1}^{\epsilon_1}(p)=p+\hat{e}_1\epsilon_1$ and $[X_1,X_j]=-\partial_{x^1} X_j =0$ implies that the components of $X_j$ do not depend on $x^1$.
Then we verify
\begin{equation}
    \phi_{X_1}^{\epsilon_1}\circ \phi_{X_j}^{\epsilon_j}(p) = \phi_{X_j}^{\epsilon_j}(p) + \epsilon_1\, \hat{e}_1= \phi_{X_j}^{\epsilon_j}(p+ \epsilon_1\, \hat{e}_1) = \phi_{X_j}^{\epsilon_j}\circ \phi_{X_1}^{\epsilon_1}(p).
\end{equation}
The rest follows by induction.
Finally for point 1 $\Rightarrow$ 3 consider the map $ \R^k\times \underbrace{M^{(n)}}_{\simeq \R^n}$ to $ \underbrace{M^{(n)}}_{\simeq \R^n}$ defined as
\begin{equation}
    F(\epsilon_1, \dots, \epsilon_k, p) := \phi_{X_1}^{\epsilon_1}\circ\dots\circ \phi_{X_k}^{\epsilon_k} (p).
\end{equation}
By point 1 the $(k+n)\times n$ jacobian matric of $F$ at any $\bar{\epsilon}\in \R^k$ and any $p\in M$ reads
\begin{equation}
    J_{F}(\bar{\epsilon}=0, p) = \begin{pmatrix}
    X_1(p) & X_2(p) & \cdots & X_k(p) & \mid &\mathbb{I}_n
    \end{pmatrix}.
\end{equation}
By the open mapping theorem, we can locally restrict $F$ to a diffeomorphism $\R^k\times \R^{n-k} \to M^(n)$ which furnishes the desired coordinates.
%
\end{proof}


\begin{mytheorem}[Flows \& Lie brackets of $f$-related vector fields]
%
For any given map $f:M\to N$ between two manifolds and vector fields $X_1,X_2\in \mathfrak{X}(M)$ and $Y_1,Y_2 \in \mathfrak{X}(N)$ that are $f$-related, i.e. such that
\begin{equation}
    f_* X_p = Y_{f(p)} \quad \forall p\in M,
\end{equation}
the flows and Lie brackets are also $f$-related, i.e.
\begin{align}
    f_* &(\phi_{X_i,\epsilon} (p)) = \phi_{Y_i,\epsilon} (f(p)) \quad \forall p\in M,
    \\[6pt]
    f_*& [X_1, X_2]_p = [Y_1, Y_2]_{f(p)} \quad \forall p\in M.
\end{align}
\end{mytheorem}


\begin{mytheorem}[Foliations, distributions \& Frobenius theorem]
%
A disribution is a sub-bundle $S \subseteq TM$ of the tangent bundle, that is a smooth assignment of a vector subspace $S_p \subseteq T_p M$ at each point $p\in M$.
A foliation is a decomposition of the manifold $M$ into a union of disjoint isomorphic connected submanifolds (the leaves) of the same dimension, so that locally the manifold looks like a product $M \simeq L^{(k)} \times N^{(n-k)}$ of a leaf $L^{(k)}$ and a `trunk' manifold $N^{(n-k)}$.
In particular for each point $p\in M$ there is a unique leaf $L_p$ passing through it, and the tangent space to the leaf at $p$ defines a subspace $T_p L_p \subseteq T_p M$.
Simple examples of foliations are
\begin{itemize}
    \item The foliation of $\mathbb{R}^3$ into planes of constant $z$ coordinate, with leaves $L^{(2)} = \{(x,y,z_0) \mid x,y\in \mathbb{R}\}$ for each $z_0\in \mathbb{R}$;
    \item The foliation of $\mathbb{R}^3\setminus \{0\}$ into spheres of constant radius, with leaves $L^{(2)} = \{(x,y,z) \mid x^2+y^2+z^2 = r_0^2\}$ for each $r_0>0$.
\end{itemize}

\textbf{Frobenius theorem}: a distribution $S \subseteq TM$ is integrable-- i.e. arises as the tangent bundle of a foliation-- if and only if it is involutive-- i.e. closed under Lie brackets.
In formulas
\begin{align}
\begin{aligned}
    &S \text{ integrable} \quad \Leftrightarrow \quad S_p = T_pL_p \,\, \forall p\in M\text{ for a foliation } M = \bigsqcup L^{(k)} 
    \quad \Leftrightarrow \quad [X,Y] \in \Gamma(S) \quad \forall X,Y \in \Gamma(S).  
\end{aligned}
\end{align}
%
\end{mytheorem}
\begin{proof}
The proof is constructive and quite long, so we just sketch the main ideas.
\begin{itemize}
    \item[$\Rightarrow$] If $S$ is integrable, then for any two vector fields $X,Y \in \Gamma(S)$-- i.e. tangent to the leaves-- their Lie bracket $[X,Y]$ is also tangent to the leaves, so that $[X,Y] \in \Gamma(S)$.
    \item[$\Leftarrow$] If $S$ is involutive, then for any point $p\in M$ we can find a local basis of vector fields $\{X_1, \dots, X_k\} \subseteq \Gamma(S)$ spanning $S$ in a neighborhood of $p$.
    By involutivity, their Lie brackets vanish $[X_i,X_j]=0$, so that by the previous theorem we can find local coordinates $\{x^1, \dots, x^n\}$ such that $X_i = \partial_{x^i}$ for $i=1\dots k$.
    Then the submanifolds defined by constant values of the last $n-k$ coordinates
    \begin{equation}
        L^{(k)}_{c_{k+1}, \dots, c_n} := \{ (x^1, \dots, x^n) \mid x^{k+1} = c_{k+1}, \dots, x^n = c_n\}
    \end{equation}
    define a foliation whose tangent spaces coincide with $S$.
\end{itemize}
\end{proof}
\todotag{Modify pic}
\begin{figure}[h]
    \centering
    \begin{tikzpicture}[scale=1.0]
        \def\R{2.1}
        \def\r{0.55}
        \pgfmathdeclarefunction{torusX}{2}{\pgfmathparse{(\R + \r*cos(#2)) * cos(#1) + 0.35*\r*sin(#2)}}
        \pgfmathdeclarefunction{torusY}{2}{\pgfmathparse{0.6*(\R + \r*cos(#2)) * sin(#1) + 0.45*\r*sin(#2)}}

        % Silhouette of the torus
        \shade[ball color=gray!25, opacity=0.8] (0,0) ellipse (2.9 and 1.8);
        \fill[white] (0,0) ellipse (1.15 and 0.6);
        \draw[line width=1.1pt, black!65] (0,0) ellipse (2.9 and 1.8);
        \draw[line width=1.1pt, black!45] (0,0) ellipse (1.15 and 0.6);

        % Leaves wrapping around the torus with an irrational slope
        \def\alpha{1.45}
        \foreach \shift/\col in {0/forestgreen!80!black, 70/forestgreen!60!black, 140/forestgreen!60!black}{
            \draw[\col, line width=1pt, samples=240, smooth, domain=0:360, variable=\u]
                plot ({torusX(\u,\alpha*\u+\shift)}, {torusY(\u,\alpha*\u+\shift)});
        }
        % A meridional leaf for contrast
        \draw[blue!70!black, line width=0.9pt, samples=200, smooth, domain=0:360, variable=\u]
            plot ({torusX(\u,\u)}, {torusY(\u,\u)});
    \end{tikzpicture}
    \caption{Sample foliation of a $2$-torus embedded in $\mathbb{R}^3$, with leaves wrapping around both directions.}
    \label{fig:torus-foliation}
\end{figure}


%-----------------------------------------------------------------------
%========================================================
\section{Derivatives on tensors}
%=========================================================
%-----------------------------------------------------------------------

There are very many notions of derivatives on tensors.
In general, a (anti)-derivative of degree $(h,k)$ on tensors is a map from a suitable subset $\mathcal{D}$ of the tensor bundle 
\begin{equation}
    \delta: T^{\bullet}_{\bullet} (M)\subseteq \mathcal{D} \to T^{\bullet + h}_{\bullet + k} (M)
\end{equation}
that is linear and satisfies a Leibniz rule, up to a sign.
That is 
\begin{equation}
    \delta (S \otimes T) = (\delta S) \otimes T + (\pm)^{pq} S \otimes (\delta T),
\end{equation}
We call it \emph{anti}derivative if it changes sign when acting on the two factors of a tensor product.
For example, the exterior derivative $d$ is an antiderivative of degree $(0,1)$ on  differential forms, and the inner derivative $i_X$ along a vector field $X$ is an antiderivative of degree $(0,-1)$ on differential forms. 


\begin{mytheorem}[Uniqueness theorem for derivatives on tensors]
An (anti)-derivative $\delta$ on tensors--i.e. a linear map and obeying Leibniz rule up to a possible minus sign-- is completely determined by its action on functions and vector fields, or equivalently on functions and 1-forms.
\end{mytheorem}



\begin{mytheorem}[The variation of a tensor is a tensor.]\todotag{move it}
Indeed the vairation of a tensor $A^\bullet_\bullet$ is simply defined as 
\begin{equation}
    \delta A^\bullet_\bullet(x):= \frac{d A^\bullet_\bullet(x,\epsilon)}{d\epsilon}_{\mid \epsilon=0} \cdot \epsilon =\lim_{h\to0} \frac{A^\bullet_\bullet(x,h)-A^\bullet_\bullet(x,0)}{h} \cdot \epsilon,
\end{equation}
where $\epsilon\to A^\bullet_\bullet(x,\epsilon)$ parametrizes a family of tensors at the same point $x\in M$, so that we can take the usual Newton derivative.
At each $|h|>0$ the ratio at the right-hand side is a tensor, being the difference and multiple of tensors, so that also the limit transforms as a tensor.
\end{mytheorem}


%========================================================
\subsection{Lie derivative}\todotag{Merge various definitions into bullet list}
\label{sec:lie_derivative}
%=========================================================

\begin{mytheorem}[Lie derivative as commutator of vector fields]
The Lie derivative $\mathcal{L}_X$ along a vector field $X$ is the unique derivative of degree $(0,0)$ on tensors that reduces to the usual directional derivative when acting on functions and to the usual Lie bracket $\mathcal{L}_X Y = [X,Y]$ when acting on vector fields.
It is then extended to all tensors using linearity and the Leibniz rule
\begin{equation}
    \mathcal{L}_X (S \otimes T) = (\mathcal{L}_X S) \otimes T + S \otimes (\mathcal{L}_X T).
\end{equation}
\end{mytheorem}


\begin{mytheorem}[Lie derivative as flow derivative]
The Lie derivative is equivalently defined as the infinitesimal generator of the flow induced by a vector field \(X\).
Namely, given a tensor field \(T\in T^{\bullet}_{\bullet} (M)\), we define the Lie derivative along \(X\) as
\begin{equation}
    \mathcal{L}_X T := \lim_{\epsilon \to 0} \frac{\phi_{X,\epsilon}^* (T) - T}{ \epsilon},
\end{equation}
where \(\phi_{X,\epsilon}: M \to M\) is the flow generated by the vector field \(X\) for parameter distance \(\epsilon\) and \(\phi_{X,\epsilon}^*\) is the pullback map induced by \(\phi_{X,\epsilon}\) on tensors.
\end{mytheorem}


\begin{mytheorem}[Cartan formula \& yet another way to define the Lie derivative]
A fundamental result linking the Lie derivative $\mathcal{L}_X$, the exterior derivative $d$ and the inner derivative $i_X$.
Namely, for any differential form $\omega$, we have
\begin{equation}
    \mathcal{L}_X \omega = i_X d \omega + d (i_X \omega).
\end{equation}
The right-hand side is defined only for differential forms, but the left-hand side makes sense for any tensor.
Since differential forms include functions (0-forms) and differentials (1-forms), we can use this formula to \emph{define} the Lie derivative and then extend it to all tensors using linearity and the Leibniz rule.
\end{mytheorem}

\begin{mytheorem}[The Lie derivative is ignorant of the curvature \& the metric]
As evident from any of its definitions, the Lie derivative does not depend on any notion of connection or metric on the manifold.
It only depends on the smooth-- i.e. differential-- structure of the manifold, being defined in terms of commutators, or flows or exterior and inner derivatives.
It thus cannot give any information on the curvature of the manifold beyond the constrained already encoded in the smooth structure itself.
\end{mytheorem}


\begin{mytheorem}[Lie derivative and pullback commute]
Given a diffeomorphism $\phi: M \to N$ between two manifolds and a vector field $Y\in \mathfrak{X}(N)$, the Lie derivative along $Y$ of any tensor field $T$ on $N$ and the pullback $\phi^*$ commute, i.e.
\begin{equation}
    \phi^* (\mathcal{L}_Y T) = \mathcal{L}_{\phi^* Y} (\phi^* T).
\end{equation}
%
More generally, if $X\in \mathfrak{X}(M)$, $Y\in \mathfrak{X}(N)$, and $S\in T^h_kM$, $T\in T^h_kN$ are $f$-related vector and tensor fields for a map $f:M\to N$, then the Lie derivatives along $X$ and $Y$ are also $f$-related, i.e.
\begin{equation}
    f_*\mathcal{L}_XS_{|p}= \mathcal{L}_Y T_{|f(p)} \quad \forall p\in M.
\end{equation}
\end{mytheorem}



\begin{mytheorem}[Lie derivative in coordinate basis \& replacing partial with covariant derivatives]
In a coordinate basis $\{\partial_\mu\}$, the Lie derivative along a vector field $X = X^{\mu} \partial_\mu$ acts on a $(h,k)$-tensor field $T$ as
\begin{align}
\begin{aligned}
    \mathcal{L}_X T^{\nu_1 \dots \nu_h}_{ \rho_1 \dots \rho_k} 
    &= X^{\mu} \partial_\mu\left( T^{\nu_1 \dots \nu_h}_{ \rho_1 \dots \rho_k} \right)
    - \sum_{i=1}^h \left(\partial_\mu X^{\nu_i}\right) T^{\nu_1 \dots \mu \dots \nu_h}_{ \rho_1 \dots \rho_k} 
    + \sum_{j=1}^k \left(\partial_{\rho_j} X^{\mu}\right) T^{\nu_1 \dots \nu_h}_{\rho_1 \dots \mu \dots \rho_k}.
\end{aligned}  
\end{align}
Note that the above expression is a tensor, even if it contains ordinary partial derivatives.
If we are given a metric $g$ \footnote{or equivalently whenever we have a torsion-less connection for which we can then find normal coordinates where the Christoffel symbols vanish} we can replace partial derivatives with covariant derivatives in normal coordinates, obtaining
\begin{align}
\begin{aligned}
    \mathcal{L}_X T^{\nu_1 \dots \nu_h}_{ \rho_1 \dots \rho_k} 
    &= X^{\mu} \nabla_\mu\left( T^{\nu_1 \dots \nu_h}_{ \rho_1 \dots \rho_k} \right)
    - \sum_{i=1}^h \left(\nabla_\mu X^{\nu_i}\right) T^{\nu_1 \dots \mu \dots \nu_h}_{ \rho_1 \dots \rho_k} 
    + \sum_{j=1}^k \left(\nabla_{\rho_j} X^{\mu}\right) T^{\nu_1 \dots \nu_h}_{\rho_1 \dots \mu \dots \rho_k}.
\end{aligned}
\end{align}
The last expression is tensorial by construction.
Since it matches the previous one in normal coordinates, it is indeed equivalent to it in any coordinate system.
\end{mytheorem}



%========================================================
\subsection{Parallel transport, connection \& covariant derivative}
%=========================================================


Parallel transport and covariant derivative are two faces of the same coin.
Namely, covariant derivative can be defined as the infinitesimal generator of parallel transport along curves.

\begin{mytheorem}[Parallel transport]
A parallel transport along a curve $\gamma: [0,1] \to M$ is a \emph{bijective linear} map 
\begin{equation}
    \tau_{\gamma}: T_{\gamma(0)} M \to T_{\gamma(1)} M
\end{equation}
such that for any two tensor fields
\begin{equation}\label{}
    \tau_{\gamma} (S \otimes T) = \tau_{\gamma} (S) \otimes \tau_{\gamma} (T).
\end{equation}
In order to obtain a well-defined notion of parallel transport, we must consistently define such a map for any curve $\gamma$ connecting any two points of the manifold.
Just like for dervatives on tensors, the parallel transport is completely determined by its action on functions and vector fields, or equivalently on functions and 1-forms.
\end{mytheorem}


\begin{mytheorem}[Connection]
A connection is a \emph{derivative} of degree $(0,1)$ on tensors, i.e. obeys linearity and Leibniz rule,
\begin{equation}
    \nabla: T^{\bullet}_{\bullet} (M) \to T^{\bullet}_{\bullet + 1} (M),
\end{equation}
that reduces to the usual directional derivative when acting on function.
It is completely determined by its action on functions and vector fields, or equivalently on functions and 1-forms.
In particular, one defines the \textbf{Christoffel symbols} $\Gamma^{\rho}_{\mu \nu}$ with respect to a coordinate basis as
\begin{equation}
    \nabla_{\partial_\mu} \partial_\nu = \Gamma^{\rho}_{\mu \nu} \partial_\rho.
\end{equation}
For the interpretation of the Christoffel symbols, see the discussion in \todotag{add ref to section}
By the Leibniz rule, the general action on a $(h,k)$-tensor field $T$ reads
\begin{equation}\label{eq:connection_on_tensor}
    \nabla_{\partial_\mu} T^{\nu_1 \dots \nu_h}_{ \rho_1 \dots \rho_k} 
    = \partial_\mu\left( T^{\nu_1 \dots \nu_h}_{ \rho_1 \dots \rho_k} \right)
    + \sum_{i=1}^h \Gamma^{\nu_i}_{\mu \sigma} \,T^{\nu_1 \dots \sigma \dots \nu_h}_{ \rho_1 \dots \rho_k} 
    - \sum_{j=1}^k \Gamma^{\sigma}_{\mu \rho_j}\, T^{\nu_1 \dots \nu_h}_{\rho_1 \dots \sigma \dots \rho_k}.
\end{equation}
Note that the value of $\nabla_X T$ at a point $p\in M$ depends only on the values of $T$ along any curve whose tangent at $p$ is $X_p$.
\end{mytheorem}

\begin{mytheorem}[Covariant derivative \& connection as infinitesimal generator of parallel transport]
Given a curve $\gamma: [0,1] \to M$ and a connection $\nabla$, the parallel transport $\tau_{\gamma}$ along $\gamma$ generated by $\nabla$ is defined as the unique map such that for any tensor field $T\in T_{\gamma(0)} M$ the curve $t \mapsto T(t) := \tau_{\gamma|_{[0,t]}} (T)$ satisfies the parallel transport equation
\begin{equation}
    0 = \nabla_{\dot{\gamma}(t)} T(t) 
\end{equation}
As usual, it is sufficient to define it on functions and vector fields, or equivalently on functions and 1-forms, and then extend it to all tensors using linearity and the Leibniz rule.
The parallel transport map is well-defined thanks to standard theorems on existence and uniqueness of solutions of ODEs.

If we are given a tensor field $S(t)$ along the curve $\gamma$, we can define the covariant derivative along $\gamma$ as
\begin{equation}
    \frac{D}{dt}S(t) := \lim_{\epsilon \to 0} \frac{\tau_{\gamma|_{[t+\epsilon,t]}}^{-1} (S(t+\epsilon)) - S(t)}{\epsilon} 
\end{equation}
That is we use parallel transport induced by $\nabla$ to compare tensors at different points along the curve.
This is indeed way $\nabla$ is called \emph{connection}: it gives a way to connect and compare `generalied tangent planes' i.e. tensors at different points.
As expected, we obtain back the covariant derivative gives back the original expression in terms of the connection
\begin{equation}
    \frac{D}{dt}S(t) := \lim_{\epsilon \to 0} \frac{\tau_{\gamma|_{[t+\epsilon,t]}}^{-1} (S(t+\epsilon)) - S(t)}{\epsilon}
    \equiv \nabla_{\dot{\gamma}(t)} S(t),
\end{equation}
where we use that \eqref{eq:connection_on_tensor} depends only on the values of $S$ along the curve whose tangent at $p$ is $\dot{\gamma}(t)$.
\end{mytheorem}



%========================================================
\subsection{The total covariant exterior derivative}
%=========================================================

See the revelant section in \ref{}




%-----------------------------------------------------------------------
%========================================================
\section{Orientation \& coverings}
%=========================================================
%-----------------------------------------------------------------------



%-----------------------------------------------------------------------
%========================================================
\section{Embeddings \& transversality}
%=========================================================
%-----------------------------------------------------------------------



%-----------------------------------------------------------------------
%========================================================
\section{Tubular neighborhoods \& collars}
%=========================================================
%-----------------------------------------------------------------------